% !TEX root = ../nc-quadric-surface-classification.tex
\section{交差項を持つ Ore 環の正規二次切断}
\label{sec:ore-normal-quadrics}

$u,v$ を中央の次数一元、$q\in k^\times$、$0\ne h\in k[u,v]_2$ とし、
\[
 A=k[u,v]\langle x,y\rangle/(yx-qxy-h)
\]
を考える。これは $k[u,v,x]$ 上の Ore 拡大であり、型 $(1,4,6,4,1)$ を持つ
（命題~\ref{prop:ore-ambient-family} と同じ Ore 拡大定理を用いる）。

\begin{proposition}
\label{prop:ore-normal-quadrics}
$A$ の非零正規二次元は、スカラー倍を除いて次の形に限る。
\begin{enumerate}
\item $q=1$ の場合は $k[u,v]_2$ の元。
\item $q\ne 1,-1$ の場合は $k[u,v]_2$ の元または $xy+h/(q-1)$。
\item $q=-1$ の場合は $a x^2+b y^2+H(u,v)$ または $xy-h/2$。
\end{enumerate}
従って $q\ne-1$ のとき、これらの切断が与える $\HH$ の標識類は、
$q\ne1$ かつ $h$ が非退化である場合の可換二次曲面一つに限る。
$q=-1$ の非中央切断も同じ条件で同じ標識類を与える。
\end{proposition}

\begin{proof}
$g$ を正規二次元とし $ag=g\nu(a)$ と書く。中央元 $u,v$ は $\nu$ で固定される。
$x,y$ の全次数を先に比較する次数フィルトレーションを取ると、
その付随次数環は $yx=qxy$ の skew polynomial 環である。
その中で $x,y$ の辞書順序をどちらから始めてもよく、
積の最高単項式では指数が加算される。
$xg=g\nu(x)$ と $yg=g\nu(y)$ の最高指数を比較すると、
$\nu(x)=\lambda x+L(u,v)$、$\nu(y)=\mu y+M(u,v)$
となり、$\lambda,\mu\ne0$ である。

$q\ne1$ なら $\nu$ を Ore 関係に適用する。
$x,y$ の一次の係数から $(1-q)\lambda M=(1-q)\mu L=0$ となり、
$L=M=0$ である。また $h\ne0$ なので $\lambda\mu=1$ である。
$g$ の $x,y$ 次数が最大の層にある単項式 $x^m y^n$ を取ると、
$x,y$ との交換から $\lambda=q^{-n}$、$\mu=q^m$ となる。
従って $q^{m-n}=1$ である。
$q\ne1,-1$ なら $m+n\le2$ と合わせて $m=n$ となる。
最大層が $xy$ の場合、$x,y$ 次数二の比較で
$u,v$ を係数とする一次の項が消える。
実際、$x$ との交換で $x$ 項は消え、$y$ との交換で $y$ 項は消える。
$g=a xy+H(u,v)$ と書くと
$xg=gq^{-1}x$ は $(q-1)H=ah$ と同値である。
最大層が次数零なら $g\in k[u,v]_2$ である。

$q=-1$ の場合、最大層が $xy$ なら同じ計算を適用する。
他の最大層では $m-n$ が偶数で、平方の層または次数零しかなく、
$\lambda=\mu=1$ となる。$x^2,y^2$ は中央元であり、
残る $x,y$ の一次の項は中心に入らないので、
中央二次式は表示した形になる。

$q=1$ なら交換子は $x,y$ 次数を下げる。
最高層の比較で $\lambda=\mu=1$ となり、
$[x,g]=gL$、$[y,g]=gM$ の次数を比較すれば $L=M=0$ である。
従って $g$ は中央である。
PBW 表示で $[y,-]$ は $h\,\partial/\partial x$、
$[x,-]$ は $-h\,\partial/\partial y$ として作用する。
標数零かつ $h\ne0$ なので中心は $k[u,v]$ である。

$k[u,v]_2$ の非零元は一次式の積に分かれ、商は domain でない。
一方 $g=xy+h/(q-1)$ で割ると $xy=yx=h/(1-q)$ となり、
商は四変数可換多項式環の二次切断である。
これは $h$ が非退化である場合にちょうど滑らかで、
その場合は可換二次曲面を与える。
元の切断環の domain 性と tails の滑らかさは、
定理~\ref{thm:sign-quadric-normal} の末尾と同じく $\HH$ への所属に必要である。
従って上の標識類以外は、この範囲では現れない。
\end{proof}

\subsection{反交換 Ore 環の中央切断}
\label{subsec:ore-central-anticommuting}

\begin{theorem}
\label{thm:ore-anticommuting-central}
$q=-1$ の中央切断 $g=a x^2+b y^2+H(u,v)$ が
$\HH$ の標識類を与えるための必要十分条件は、
$a,b$ がともに非零であり、以下の二つの二次形式 $f_1,f_2$ が非退化となることである。
$x,y$ をスカラー変更して $a=b=1$ とし、その規約で $f_1=-H+h$、$f_2=-H-h$ とする。
許容される標識類は、$\PP^1$ 上の相異なる二点の非順序対を二つ取り、
対の交換と $\PGL_2$ で割ったデータで判定される
（二つの対は互いに一致してもよい）。どの類も標識付き cubic 表示を持たない。
\end{theorem}

\begin{proof}
$a=b=0$ なら商は domain でない。
一方だけ零、例えば $b=0$ なら $B^!$ で $w=X^2$ は中央正則で、
$Y^2=0$、$XY=YX$ である。$X,Y$ は $u,v$ の双対生成元と反交換する。
従って $C(B)$ 中で $YX^{-1}$ は非零中央平方零元であり、
この場合は滑らかでない。

次に $a=b=1$ とする。$c=x+y,d=x-y$ と置けば
\[
 B=k[u,v]\langle c,d\rangle/(c^2-f_1,d^2-f_2).
\]
交互に並ぶ $c,d$ の語が $k[u,v]$ 上の基底である。
書換え $c^2\mapsto f_1$、$d^2\mapsto f_2$ は長さを下げ、
三文字の同じ文字の重なりは係数の中心性で解消するからである。
元の Ore 環の関係は $c^2-d^2=f_1-f_2$ である。
双対で $CD=DC=0$、$w=C^2+D^2$ は正則正規二次元の代表である。
$C^2,D^2$ は中央なので、$C(B)$ は直交冪等元 $C^2/w,D^2/w$ により分解する。
第一成分では $C$ が対応する局所化の角で可逆となり、
$UC^{-1},VC^{-1}$ が二次形式が $f_1$ に比例する二元 Clifford 環を生成する。
第二成分では $D$ と $f_2$ を用いた同じ計算になる。
実際、$C,D$ は $U,V$ と反交換し、
$U,V$ の平方と反交換子は $f_1,f_2$ の係数を掛けた $C^2,D^2$ の和である。
従って各成分は高々四次元で、$\dim C(B)=8$ により両方が四次元である。
標数零なので二元 Clifford 環は非退化の場合にちょうど半単純であり、
退化の場合には根基のベクトルが非零平方零イデアルを生成する。
従って滑らかさは $f_1,f_2$ がともに非退化であることと同値である
\cite[Lemma 4.13 and Theorem 5.5]{MU2022}。

以後 $f_1,f_2$ を非退化とする。$B$ は domain である。
$F=k(u,v)$ に移し、中央元 $z=cd+dc$ を取ると、
交互語の基底から $F[z]$ 上で $1,c,d,cd$ が自由基底になる
（$z$ の冪の最高長さの交互語を比較する）。
さらに $F(z)$ に移し $e=d-zc/(2f_1)$ とすれば、
$c^2=f_1$、$ec=-ce$、$e^2=-(z^2-4f_1 f_2)/(4f_1)$ の四元数環となる。
$f_1,f_2$ はともに $F$ の平方でない。
$z^2-4f_1 f_2$ の任意の既約因子の素点では、$f_1$ は剰余体でも平方ではない。
$f_1 f_2$ が平方なら剰余体は $F$ であり、そうでなければ $F(\sqrt{f_1 f_2})$ である。
後者で $f_1$ が平方なら $F$ で $f_1$ または $f_2$ が平方になってしまう。
従ってこの素点における $\sqrt{f_1}$ の拡大は非分岐で剰余次数二となる。
ノルムの付値は偶数だが $e^2$ の付値は一である。
よって $e^2$ はノルムでなく、この四元数環は斜体である。
自由性により元の $B$ も domain となる。

適合根と被覆を構成する。
$f_1,f_2$ の二つの根の対を互いに交換する $\PP^1$ の対合を取れる。
根の対が一致するときは恒等を使う。
一根を共有するとき、座標で二対を $\{0,\infty\},\{0,1\}$ とすれば
$t\mapsto t/(t-1)$ が働く。四根が相異なるときは
$\{0,\infty\},\{1,\lambda\}$ とし、$t\mapsto\lambda(t-1)/(t-\lambda)$ を使う。
二乗が恒等の線形持上げ $L$ を選び、
$Lf_1=\eta f_2,Lf_2=\eta^{-1}f_1$ とする。
$\alpha^2=\eta$ を選び、$\sigma$ を $u,v$ 上で $L$、
$c\mapsto\alpha d$、$d\mapsto-\alpha^{-1}c$ とする。
$r_c=c^2-f_1,r_d=d^2-f_2$ と書けば $\sigma$ は
$r_c-\eta r_d$ の直線と $r_c+\eta r_d$ の直線を保つ。
$A'$ を中央 $u,v$ と第一の関係 $r_c-\eta r_d$ だけで定め、
第二の元 $g'=r_c+\eta r_d$ で切断すれば $B$ が得られる。
$c\pm\sqrt\eta\,d$ を新しい変数に取ると、
$A'$ は反交換 Ore 環である。型 $(1,4,6,4,1)$ を持ち、
$c^2,d^2$ が中央であるため $g'$ は中央正則である。

$A'$ の Nakayama 自己同型は $u,v$ を固定し $c,d$ に負号を掛ける。
Ore 表示の基礎環での作用は \cite[Theorem 3.3]{LWW2014} により、
残る変数の作用は Ore 関係との整合から従う。
Ore 関係の $k[u,v]$ 項が非零なら、基礎環の非中央変数に関する次数を比較して
Ore 関係を保つために平行移動項は零となり、係数は $-1$ である。
Ore 関係の $k[u,v]$ 項が零なら符号型の Nakayama 公式を使える。
従って中央切断の公式 \cite[Lemma 1.5]{RRZ2014} により、$B$ の Nakayama 自己同型も
$u,v$ を固定し $c,d$ を反転する。これは $\sigma^2$ と一致する。
よって $B^\sigma=A'^\sigma/(g')$ の通常シフトが適合根となる。
飽和性、例外性、豊富性は前の符号型の場合と同じ正則切断の議論による。

GK 条件には、$B$ が $k[u,v]$ 上自由で、
$B/(u,v)=k\langle c,d\rangle/(c^2,d^2)$ が無限次元であることを使う。
$B/uB$ 上の $v$ の冪の列は tails で真に下降するので、
$B/uB$ の像は有限長でない。
$u$ の冪の列と domain 性を使うと、前と同様に rank one torsion-free 性が従う。

$f_1=-\ell_1\ell_2$ と一次式の積に書けば、
行列 $\left(\begin{smallmatrix}c&\ell_1\\ \ell_2&-c\end{smallmatrix}\right)$ の二乗は
$r_c I$ であり、$A'$ 上で $g'$ の非零スカラー倍になる。
$f_2$ と $d$ に対しても同様である。
二型の成分空間はそれぞれ $\langle u,v,c\rangle$ と
$\langle u,v,d\rangle$ で、両方が階数三である。
twist は階数を保つので標識付き cubic 表示は存在しない。

最後に二つの許容切断を比較する。
適合環の添字付き代数は、それぞれ twist 前の $B$ の添字付き代数と
同型である。従って命題~\ref{prop:first-orbit-finite-equivalence} を使える。
上の $B$ の双線形化関係は、$(u,v)\ne(0,0)$ では
\[
 E=\{c^2=f_1(u,v),\ d^2=f_2(u,v)\}\subset\PP^3
\]
の上の対角対応を与える。$u,v$ のいずれかを使う交換関係から
次の点が比例するためである。$(u,v)=(0,0)$ では、
$c,d$ の二つの座標点が互いに移る。これらは $E$ には含まれない。
二つの多項式の根の対が異なれば $E$ は整で $\PP^3$ を張る。
実際、$\PP^1$ の関数体上で独立な二平方根を付加する被覆であり、
式は基底 $1,c,d,cd$ の有限自由被覆を与える。
根の対が一致すれば $E$ は異なる二平面にある二つの滑らかな円錐曲線であり、
二点で交わる。いずれの場合も $E$ を点ごとに固定する射影線形変換は恒等である。
整の場合は固有空間のいずれか一つに全体が入ることから、
二成分の場合は二平面上のスカラーが交点で一致することから従う。

従って添字付き同型の連続する射影線形変換は $E$ 上で一致するので、
スカラー倍まで同じである。よって同値は $B$ の次数付き同型に帰着する。
次数一の中心は $\langle u,v\rangle$ である。
$c$ の像を $a c'+b d'+\ell(u,v)$ と書くと、平方が中央変数の多項式であるため
交互語の基底から $ab=0,\ell=0$ を得る。$d$ に対しても同様である。
従って次数付き同型は $u,v$ の基底変更と、$c,d$ の交換およびスカラー変更
に尽きる。これはちょうど二つの根の対に対する主張した同値関係である。
逆に任意の二対を定める非退化な二次形式を取り、
必要なら一方をスカラー倍して $f_1-f_2\ne0$ とすれば、この Ore 族から実現できる。
\end{proof}

\begin{corollary}
多パラメータ skew polynomial の切断と、この節の Ore 環の切断を合わせると、
$\HH$ に属する標識類は、$r^2$ を逆元で割った cubic 族と、
定理~\ref{thm:ore-anticommuting-central} の非 cubic 族である。
後者は根の対が一致する一点、根を一つ共有する一点、
および四根が異なる場合の $\lambda\in k\setminus\{0,1\}$ を
$\lambda\leftrightarrow\lambda^{-1}$ で割った族からなる。
\end{corollary}
\begin{proof}
前者は定理~\ref{thm:multiparameter-quadric-classification} による。
同定の際、多パラメータ切断の非 cubic の一点は
命題~\ref{prop:nonstandard-rank-three-example} を $c,d$ の表示で書けば
両平方が同じ二次形式となる場合である。
後者の四根が異なる場合は二対を $\{0,\infty\},\{1,\lambda\}$ へ移せる。
各対内の交換は交比を逆にし、二対そのものの交換は交比を保存する。
一致と一根共有の場合は、それぞれ二点または三点に対する
$\PGL_2$ の推移性から一点となる。
\end{proof}
